\chapter{Conclusion and Future Work}
\label{chap:conclusion}
\section{Conclusion}
In the single-step Gaussian fusion experiments, this thesis finds that NFF exhibits clear advantages in several typical challenging scenarios. In particular, when the likelihood distribution is very narrow, when the prior and the likelihood are far apart, or when the system dimension increases continuously, PF and GPF are more prone to significant estimation errors, whereas NFF is still able to maintain relatively stable estimation results. The experiments also show that the choice of NFF parameters affects its performance. In particular, there exists a suitable range for the particle count, and a larger number of particles does not necessarily lead to better results. Similarly, under the experimental conditions considered in this thesis, the covariance-constrained version does not improve performance. In addition, the estimation accuracy of NFF is not very sensitive to the shape of the homotopy path, but a suitable path can reduce the stiffness of the ODE and thereby decrease the computational time. Under multimodal distributions, NFF tends to track only one of the modes.

In the Lorenz 96 experiments, when the observation noise is very small, NFF clearly outperforms PF and GPF because it is better able to avoid particle degeneracy. As the observation noise increases, the performance gap among the three methods gradually narrows. When the system process noise increases, the error of NFF grows more slowly, indicating that it is better suited to complex systems with strong noise. When the system nonlinearity becomes stronger, NFF can still exploit the prior information more effectively than PF and GPF, as long as the prediction step is still able to provide some useful information. However, if the nonlinearity becomes so strong that the prediction itself fails, the performance of NFF also deteriorates significantly. The model mismatch and extreme-condition experiments further show that NFF has good stability and robustness in most cases.

Compared with PF and GPF, NFF generally achieves better estimation performance under high-precision observations, large process noise, strong nonlinearity, and certain types of model mismatch, while not requiring the particle count to be increased as drastically as in PF. Moreover, the third prediction mechanism of NFF performs significantly better than the first prediction mechanism. However, the computational time of NFF is much longer than that of PF and GPF.

\section{Future Work}

Although this thesis has systematically compared the PF, the GPF, and the NFF on the fully observed Lorenz 96 system, several important directions remain for future investigation.

First, the present study considers only the fully observed Lorenz 96 system. In practice, partially observed Lorenz 96 systems are often more challenging, especially for PF and GPF. Therefore, it is an important next step to investigate whether NFF can still maintain its advantages under partially observed settings.

Second, additional comparisons with other filtering methods would further clarify the strengths and limitations of NFF. Since the Ensemble Kalman Filter (EnKF) is widely regarded as particularly suitable for the Lorenz 96 system, a comparison between EnKF and NFF would be highly valuable. Moreover, the observation model considered in this thesis is linear, and the Unscented Kalman Filter (UKF) is also well suited to such settings. Therefore, comparing NFF with UKF would provide a more comprehensive evaluation framework.

Third, the computational efficiency of NFF can be further improved. The current NFF is based on the CVM distance defined through the LCD with a Gaussian kernel. Since the Gaussian kernel has infinite support, each particle is coupled with all other particles, which leads to a high computational cost. A promising direction for future work is therefore to reformulate the LCD and the corresponding CVM distance such that each particle interacts only with its nearest neighboring particles. This may reduce the computational complexity.